\documentclass[10pt,a4paper]{amsart}
\usepackage[margin=19mm]{geometry}
\usepackage{amsmath,amssymb,mathtools}
\usepackage[scr=rsfs,cal=euler]{mathalfa}
\usepackage{xfrac}
\usepackage[unicode,hidelinks,bookmarks=false]{hyperref}
\newcommand{\ie}{i.e.\ }
\newcommand{\cf}{cf.\ }
\newcommand{\resp}{respectively\ }
\newcommand{\Subsection}{\subsection}
\providecommand{\nobreakdash}{\nobreak\hbox{-}}
\setlength{\emergencystretch}{2em}
\multlinegap0pt
\usepackage{amssymb}
\usepackage{enumerate}
\usepackage[shortlabels]{enumitem}
\usepackage[normalem]{ulem}
\newtheorem{lemma}{Lemma}
\newcommand\bN{{\mathbb N}}
\newcommand\bR{{\mathbb R}}
\newcommand\cF{{\mathcal F}}
\newcommand\cG{{\mathcal G}}
\newcommand\cI{{\mathcal I}}
\newcommand\cS{{\mathcal S}}
\newcommand\cW{{\mathcal W}}
\newcommand\cZ{{\mathcal Z}}
\newcommand{\catt}{\cF_{A,t}}
\renewcommand{\ge}{\geqslant}
\renewcommand{\le}{\leqslant}
\title{On linear response for discontinuous perturbations of smooth endomorphisms}
\author{Giovanni Canestrari}
\date{Source: arXiv:2411.16628v3 (CC BY 4.0)}
\begin{document}
\maketitle

\noindent\emph{Source and licence.} On linear response for discontinuous perturbations of smooth endomorphisms. Version arXiv:2411.16628v3 (CC BY 4.0), distributed by arXiv under the Creative Commons Attribution 4.0 International licence, http://creativecommons.org/licenses/by/4.0/, as recorded in the arXiv licence metadata for this exact version and shown on the abstract page https://arxiv.org/abs/2411.16628v3. Full licence notice: \texttt{licenses/arxiv-2411.16628-CC-BY-4.0.txt}.\par\smallskip

\noindent\textit{Editorial scope.} Counted unit: the article's lemma lem:growth-lemma (source0.tex lines 1645--1657). The article's complete original proof of it is delivered with it (source0.tex lines 1658--1690, one \texttt{proof} environment of 33 lines). No mathematics is written, repaired or re-derived: the statement and the proof are the article's own lines, in the article's own notation, and no retained line is substituted or re-flowed.  Retained uncounted as the source-local context the proof needs: lines 1575--1580; lines 1585--1599; lines 1633--1635.  Nothing was omitted from any retained span, so the package carries no omission record.  No retained line contains a citation, so the wrapper carries no bibliography and no bibliography entry.  No retained line contains a figure environment, an image-inclusion command or drawing code, so no image is delivered and the package declares no asset dependency.  Editorial formatting only, in addition: an AMS \texttt{amsart} preamble that restates the article's own declarations and macros used by the retained lines, namely the environments \texttt{lemma} on separate counters, together with the standard packages those lines need.  The retained environments print in this document as Lemma 1 in order of appearance, whereas the article prints the same items with the numbering of its own full document.  The complete native source of the article as distributed in the arXiv e-print for this exact version is bundled unchanged as \texttt{sources/169081/source0.tex}.
\begin{lemma}\label{lem:hyp-beats-cutting}
  There exists \(N_0 \in \bN\) such that,
  \[
       \frac{K_{N_0} +1}{\lambda^{N_0}} < 1.
  \]
\end{lemma}

\begin{lemma}\label{lem:trasversality2}
  For any \(k \in \bN_0\), there exists \(P_k \in \bR^+\) such that, for any \(t \in [0,1/8]\), \(W \in \cW\),
  \[
  \begin{split}
    \#\{W \cap \cS_t^{k,+}\} \le K_{k} + P_{k}|W|,
  \end{split}
  \]
  or, equivalently,
  \[
  \begin{split}
\# \{W': W' \text{ is a maximal connected}& \text{ component} \\
&\text{ of }W \setminus \cS_t^{k,+}\} \le K_{k}+1 + P_k |W|.
  \end{split}
  \]
\end{lemma}

\begin{lemma}\label{lem:weigths-single-curve}
 Let \(W \in \cW\), \(n \in \bN_0\) and denote (renaming the indices) \(\catt^n \cG_{W} = \left\{(W_q, p_q)\right\}\). Then \(p_q = |W_q|/|\catt^n (W)|\) and in particular \(p_q \ge  |W_q|/(\overline \lambda^n |W|)\).
\end{lemma} 

\begin{lemma}[Growth Lemma]\label{lem:growth-lemma}
  There exist \(z \in (0,1)\) and \(Z \in \bR^+\) such that, for any \(t \in [0, 1/8]\), \(k \in \bN_0\), and standard family \(\cG\),
  \[
      \cZ(\catt^{k N_0} \cG) \le z^k \cZ (\cG) + Z.
  \]
  Furthermore, there exist \(Z_{*}, Z_{**}, Q_1,Q_2 \in \bR^+\), such that, for all \(t \in [0, 1/8]\),
  \[
  \begin{split}
    &\cZ(\catt^{p}\cG) \le \max\{Z_{**}\cZ(\cG), Z_{*}\}, \quad \forall p \in \bN_0,\\
    &\cZ(\catt^{p}\cG) \le Z_{*}, \quad \forall p \ge Q_1\left|\ln \cZ(\cG)\right| + Q_2.
  \end{split}
  \]
\end{lemma}

\begin{proof}
 Let \(\cG = \{W_j, p_j\}_{j \in \cI}\). For each \(j \in \cI\), let \(\{W_{j,j'}\}\), \(j' \in \cI_j\), be a partition of \(\catt^{p}(W_j)\), \(p \in \bN_0\), in connected components. The cardinality of \(\cI_j\) is not bigger than the number of connected components of \(W_j \setminus \cS^{+,p}_t\). Therefore, by Lemma \ref{lem:trasversality2}, there exists a constant \(P_{p} \in \bR^+\), such that, for each \(j \in \cI\),
\begin{equation}\label{eq:cardinality-new-partition}
  \# \cI_j \le K_{p} + 1 + P_{p}|W_j|.
\end{equation}
Moreover, by Lemma \ref{lem:weigths-single-curve}, the weight associated to \(W_{j,j'}\) is
\begin{equation}\label{eq:iterated-evol}
p_{j,j'} = p_j \frac{|W_{j,j'}|}{|\catt^p(W_j)|}.
\end{equation}
Hence, by \eqref{eq:cardinality-new-partition} and \eqref{eq:iterated-evol},
\begin{equation}\label{eq:estimate-evol-Z}
\begin{split}
  \cZ(\catt^{p}\cG) &= \sum_{j \in \cI}\sum_{j' \in \cI_{j}} p_{j,j'} \frac{1}{|W_{j,j'}|}= \sum_{j \in \cI}\sum_{j' \in \cI_{j}}p_j\frac{|W_{j,j'}|}{|\catt^{p} W_j|}\frac{1}{|W_{j,j'}|} \\
  &\le \sum_{j \in \cI}p_j\frac{K_{p} +1 +P_{p}|W_j|}{\lambda^{p}|W_j|}\le \frac{K_{p} + 1}{\lambda^{p}}\cZ(\cG) + \frac{P_{p}}{\lambda^{p}}.
\end{split}
\end{equation}
We now prove the first statement of the Lemma. The case \(k = 1\) follows by \eqref{eq:estimate-evol-Z} and Lemma \ref{lem:hyp-beats-cutting}. For general \(k \in \bN\), it follows by induction, using \eqref{eq:estimate-evol-Z} with  \(\catt^{k N_0} \cG\) in place of \(\cG\) and \(p = N_0\). Moreover, equation \eqref{eq:estimate-evol-Z} with \(p = 1\) yields
\begin{equation}\label{eq:bad-estimate}
  \cZ \left(\catt \cG\right) \le\frac{K_{1} + 1}{\lambda}\cZ(\cG) + \frac{P_{1}}{\lambda} \le\frac{K_1 + 1+ \sqrt{2} P_1}{\lambda} \cZ(\cG),
\end{equation}
where we have used that \(\cZ(\cG) \ge 1/\sqrt{2}\) (recall the upper bound for the lengths of standard segments). Writing \(p = kN_0 + q\), \(k \in \bN_0\), \(q \in \{0,...,N_0-1\}\),
\begin{equation}\label{eq:growth-eq-proof}
\begin{split}
\cZ(\catt^p \cG) &\le \left( \frac{K_1 + 1+ \sqrt{2}P_1}{\lambda}\right)^q \cZ(\catt^{kN_0}\cG) \le Z_{**} \max\left\{\cZ(\cG), \frac{Z}{1-z}\right\} \\
&\le \max\{Z_{**}\cZ(\cG), Z_{*}\},
\end{split}
\end{equation}
where we have used \eqref{eq:bad-estimate} repeatedly in the first inequality, the first part of the Lemma in the second, and we have set
\begin{equation}\label{eq:star-2star}
Z_{**} = \left(\frac{K_1 + 1+ \sqrt{2}P_1}{\lambda} + 1\right)^{N_0}, \quad Z_{*} = \frac{2ZZ_{**}}{1-z}.
\end{equation}
The third claim can be obtained by the first part of the Lemma together with the first inequality in \eqref{eq:growth-eq-proof} and the definition of \(Z_{*}\) in the equation above.
\end{proof}

\end{document}
